\chapter{The Research Boundary Object}\label{ch04:boundary}

\section{The problem the boundary object answers}\label{ch04:problem}

Chapter 1 identified unauthorized claim promotion as one of the two failure
classes that motivate the kernel. The failure has a characteristic shape: an
artifact that is merely \emph{consistent}---a green test run, a self-computed
hash, a well-formed receipt, a ledger row containing the word
\textsc{licensed}---is read by some downstream consumer as \emph{authorized},
and a research candidate crosses into product territory without any reviewer,
authority, or legal act having occurred. The program's own history contains the
canonical instance: AI-authored, self-hashed ``licensing receipts'' that were
structurally impeccable and epistemically worthless. Nothing in those artifacts
was false in the narrow sense; the failure was that a status which can only be
conferred by an external act was inferred from internal structure.

The research boundary object is the kernel's answer to this failure class at
the level of representation. Its design premise is that the status of an
artifact---where it sits in the admission process, whether it may back a
product claim, whether it may be publicly disclosed---must be (i) explicit
machine-visible data rather than convention, (ii) constrained so that the
historically observed contradictory combinations cannot be constructed at all,
and (iii) rechecked at every point of use rather than trusted from any earlier
point of construction. The present chapter develops each of these three
commitments, then dwells at length on the distinction that makes the whole
design coherent: the difference between \emph{tamper evidence} and
\emph{authority}. The chapter closes with worked scenarios of boundary misuse
and the exact mechanism that blocks each one.

Throughout, the claim discipline of the monograph applies. Everything
described here is a Ring-4 research candidate: admission \textsc{open}, product
\textsc{not licensed}, disclosure withheld. The chapter describes what the
v0.1 candidate's code enforces under its bounded test fixtures; it does not
claim security against arbitrary in-process adversaries, does not claim that
any boundary object confers or withholds legal status, and does not promote
any artifact past the boundary it carries.

\begin{boundary}[Scope of this chapter]\label{ch04:bd:scope}
The boundary object is a representational and procedural discipline inside a
single Python process. It is not an access-control system, not a signature
scheme, and not a legal instrument. Statements of the form ``the kernel
rejects $X$'' mean that the v0.1 candidate raises a typed validation error on
the fixture family exercising $X$ in the bound Gate-4 suite; they are
historical software receipts, not guarantees over untested inputs.
\end{boundary}

\section{The boundary triple and its carrier}\label{ch04:triple}

\subsection{The triple}

The governing specification defines the research boundary as a triple.

\begin{definition}[Research boundary triple]\label{ch04:def:triple}
A \emph{research boundary} is the triple
\[
(\text{admission state},\ \text{product state},\ \text{public disclosure}),
\]
where the admission state ranges over $\{\textsc{open}, \textsc{in-review},
\textsc{accept}\}$, the product state over $\{\textsc{not licensed},
\textsc{licensed}\}$, and public disclosure over $\{\mathrm{false},
\mathrm{true}\}$.
\end{definition}

The three components answer three different questions, and the design treats
the questions as genuinely independent axes that are then coupled by explicit
consistency rules rather than being collapsed into a single status enum.

The \emph{admission state} answers: where does this artifact sit in the
review process? \textsc{open} means no review act has been recorded;
\textsc{in-review} means a review is underway; \textsc{accept} means an
admission decision has been recorded. Nothing about the admission state says
what the artifact may be used for.

The \emph{product state} answers: may this artifact back a product claim?
The only value that permits it is \textsc{licensed}, and---as the consistency
rules below make precise---that value is constructible only in a configuration
the v0.1 kernel never occupies. \textsc{not licensed} is not a pejorative; it
is the honest default of every research artifact.

The \emph{public disclosure} flag answers: has disclosure of this artifact
been authorized? The flag is deliberately boolean rather than an enum of
disclosure channels: v0.1 needs only the distinction between ``no disclosure
authorization exists'' and ``a disclosure authorization has been recorded,''
and a richer disclosure model would require exactly the authenticated
trust root that v0.1 deliberately lacks.

\subsection{The carrier and the ring index}

In the v0.1 implementation the triple is carried by a frozen record with four
fields, because the consistency rules are indexed by a fourth datum: the
admission ring.

\begin{definition}[Boundary carrier]\label{ch04:def:carrier}
The boundary carrier is the immutable record
\[
(\texttt{admission\_ring},\ \texttt{status},\ \texttt{license\_status},\
\texttt{public\_disclosure\_authorized}),
\]
where \texttt{admission\_ring} is an integer in $\{1,2,3,4\}$,
\texttt{status} is an admission-state enumeration value,
\texttt{license\_status} is a product-state enumeration value, and
\texttt{public\_disclosure\_authorized} is a boolean. Each field is validated
for exact runtime type at construction: the ring must be exactly an integer
(a boolean, though an integer subtype in the host language, is rejected by an
exact-type check elsewhere in the kernel and by the enumeration checks here),
the two states must be exactly members of their enumerations---not strings
that happen to spell the same word---and the disclosure flag must be exactly
a boolean.
\end{definition}

The ring is the program's admission-severity index inherited from the Project
Coherence process: Ring 1 and Ring 2 are reserved for reviewed mathematical
material, Ring 3 for artifacts that have passed an admission decision, and
Ring 4 for open research candidates. The kernel itself assigns no meaning to
the rings beyond the consistency constraints below; in particular, nothing in
the kernel can move an artifact between rings, because the carrier is
immutable and every v0.1 construction site requires the exact Ring-4
configuration.

Two further construction rules deserve emphasis before the consistency rules,
because they recur throughout the kernel and were motivated in Chapter 3.

\begin{ruleenv}[Exact carrier type]\label{ch04:rule:exacttype}
A subclass of the boundary carrier is rejected at construction: the
constructor checks that the runtime type of the object under construction is
exactly the carrier type, not a derived type. Consequently no object with
overridden equality, overridden field access, or additional behavior can pose
as a boundary value at any kernel interface.
\end{ruleenv}

\begin{ruleenv}[Exact-field predicate]\label{ch04:rule:predicate}
The kernel's boundary predicate---``is this the exact Ring-4 open
boundary?''---does not call any overridable equality. It checks, field by
field, that the value has exactly the carrier type, that the ring is exactly
the integer 4, that the status is exactly the \textsc{open} enumeration
member (by identity, not by string comparison), that the product state is
exactly \textsc{not licensed}, and that the disclosure flag is exactly the
boolean false.
\end{ruleenv}

Rule~\ref{ch04:rule:predicate} may look paranoid in isolation. Its point is
architectural: every consumer of a boundary decides for itself, from
primitive fields, what the boundary says. There is no delegation to a method
the value itself supplies, because a value that answers questions about its
own trustworthiness is exactly the anti-pattern the kernel exists to
eliminate. This is the first appearance of a theme that
Section~\ref{ch04:recheck} makes into an invariant.

\section{Joint consistency and construction-time rejection}\label{ch04:consistency}

\subsection{The three consistency rules}

The four fields of the carrier are not independent. Three rules couple them,
and the constructor enforces all three, so that a violating combination is
never a value at all---it is a raised validation error.

\begin{ruleenv}[Ring-4 closure]\label{ch04:rule:c1}
If \texttt{admission\_ring} $=4$, then the status must be \textsc{open}, the
product state must be \textsc{not licensed}, and the disclosure flag must be
false. Equivalently: a Ring-4 artifact is open, unlicensed, and undisclosed,
with no exceptions representable.
\end{ruleenv}

\begin{ruleenv}[License precondition]\label{ch04:rule:c2}
If the product state is \textsc{licensed}, then \texttt{admission\_ring}
$=3$ and the status is \textsc{accept}. Equivalently: a licensed product
state is representable only for an artifact that has a recorded admission
decision at Ring 3.
\end{ruleenv}

\begin{ruleenv}[Disclosure precondition]\label{ch04:rule:c3}
If the disclosure flag is true, then the status is \textsc{accept}.
Equivalently: no disclosure authorization is representable for an artifact
that has not passed admission.
\end{ruleenv}

Each rule encodes one historically observed contradiction as an
unrepresentable state. Rule~\ref{ch04:rule:c1} makes the phrase ``a licensed
Ring-4 candidate'' a type error rather than a category the program has to
argue about. Rule~\ref{ch04:rule:c2} makes ``licensed but never admitted''
unrepresentable---this is precisely the shape of the historical self-issued
licensing receipts, which asserted a product status with no admission act
behind it. Rule~\ref{ch04:rule:c3} makes ``disclosed but still open''
unrepresentable, which blocks the quieter failure in which a draft escapes
into public circulation and the disclosure is then rationalized after the
fact.

It bears stating plainly what these rules do \emph{not} say.
Rule~\ref{ch04:rule:c2} does not say that Ring 3 plus \textsc{accept}
\emph{makes} an artifact licensed; it says only that outside that
configuration a licensed product state cannot even be written down. Whether
a Ring-3 accepted artifact is actually licensed is a question about acts in
the world---reviewer identity, legal authority---about which the carrier
carries no evidence whatsoever. The consistency rules are necessary-condition
filters, never sufficient-condition grants. This asymmetry is the entire
content of Section~\ref{ch04:tamper}, and it is why the kernel's runtime
gate (Chapter 9) denies claim-bearing intent even when handed a
structurally valid \textsc{licensed} boundary.

\subsection{Enumerating the state space}

Because the component sets are small, the effect of the three rules can be
stated exactly.

\begin{proposition}[Valid boundary count]\label{ch04:prop:count}
The unconstrained carrier state space has $4 \times 3 \times 2 \times 2 = 48$
combinations. Exactly $15$ of them satisfy
Rules~\ref{ch04:rule:c1}--\ref{ch04:rule:c3}; the remaining $33$ are rejected
at construction.
\end{proposition}

\begin{proof}[Case analysis]
Consider each ring. For ring $4$, Rule~\ref{ch04:rule:c1} admits exactly one
of the twelve combinations: (\textsc{open}, \textsc{not licensed}, false).
For rings $1$ and $2$, Rule~\ref{ch04:rule:c2} excludes the product state
\textsc{licensed} entirely (six combinations per ring); among the six
remaining unlicensed combinations, Rule~\ref{ch04:rule:c3} excludes a true
disclosure flag for \textsc{open} and \textsc{in-review} (two combinations
per ring), leaving four per ring, hence eight. For ring $3$, the unlicensed
half contributes the same four; the licensed half is admitted by
Rule~\ref{ch04:rule:c2} only with status \textsc{accept}, for which
Rule~\ref{ch04:rule:c3} permits either disclosure value, contributing two.
Hence $1 + 8 + 6 = 15$.
\end{proof}

Table~\ref{ch04:tab:invalid} organizes the thirty-three rejected
combinations into families by the rule that rejects them; a combination
violating more than one rule is listed under the first rule that fires in
the constructor's check order.

\begin{table}[ht]
\centering
\begin{tabular}{llr}
\toprule
Family & Description & Count \\
\midrule
R4/status & ring 4 with status \textsc{in-review} or \textsc{accept} & 16 \\
R4/license & ring 4, status \textsc{open}, product \textsc{licensed} & 2 \\
R4/disclosure & ring 4, otherwise closed, disclosure true & 1 \\
Lic/ring & \textsc{licensed} at ring 1 or 2 & 8 \\
Lic/status & \textsc{licensed} at ring 3 without \textsc{accept} & 4 \\
Disc/status & disclosure true without \textsc{accept} (rings 1--3, unlicensed) & 2 \\
\bottomrule
\end{tabular}
\caption{The 33 invalid carrier combinations by rejecting family. Counts
reflect the constructor's check order, which tests the Ring-4 closure first;
the R4/status family therefore absorbs ring-4 combinations that also violate
the license or disclosure preconditions.}
\label{ch04:tab:invalid}
\end{table}

The counting exercise is not an ornament. It is the concrete meaning of the
design slogan ``invalid states are unrepresentable'': the program's
historically dangerous statuses correspond to identifiable cells of a finite
table, and each such cell is a constructor exception, testable one fixture at
a time. The Gate-4 suite exercises representatives of each family; the
receipt for that suite is bounded software evidence in the sense of
Chapter 12, no more.

\subsection{The distinguished constructor}

The kernel provides a single distinguished factory: the Ring-4 open boundary,
producing the unique valid ring-4 carrier
(\textsc{open}, \textsc{not licensed}, false). Every artifact type in the
v0.1 kernel---semantic crosswalks, registry records, registries,
compatibility receipts---defaults its boundary field to a \emph{fresh}
instance from this factory. Two properties of this arrangement matter.

First, freshness: the default is a factory invoked per construction, not a
module-level shared instance. No two artifacts share a boundary object, so no
single mutation---accidental or otherwise---can silently change the recorded
status of a population of artifacts. Sharing a mutable status object across
artifacts would create precisely the ``shared mutable trust'' that
Invariant~\ref{ch04:inv:recheck} forbids.

Second, bias: the distinguished factory is the \emph{most restrictive}
boundary, and it is the only one with a privileged constructor. Constructing
any other valid boundary---say, a Ring-3 accepted licensed one---requires
spelling out every field explicitly. The kernel does not make the dangerous
configurations impossible to write down when they are internally consistent
(Rule~\ref{ch04:rule:c2} deliberately leaves the Ring-3 licensed shape
representable, because a future reviewed design will need it); it makes them
impossible to \emph{stumble into}, and---as the next section shows---
impossible to \emph{use} anywhere in v0.1.

\subsection{Formal summary}

It is convenient to fix notation used in the rest of the chapter. Let
\[
\mathcal{B} \;=\; R \times S \times L \times D, \qquad
R = \{1,2,3,4\},\ |S| = 3,\ |L| = 2,\ |D| = 2,
\]
be the unconstrained carrier space, and let
$\mathcal{V} \subseteq \mathcal{B}$ be the subset satisfying
Rules~\ref{ch04:rule:c1}--\ref{ch04:rule:c3}, so that
$|\mathcal{B}| = 48$ and, by Proposition~\ref{ch04:prop:count},
$|\mathcal{V}| = 15$. The constructor computes the characteristic
function $\chi_{\mathcal{V}}$ and raises on the complement; the use-site
predicate of Rule~\ref{ch04:rule:predicate} computes the characteristic
function of a much smaller set---the singleton containing the Ring-4
element---together with the exact-type conditions that
$\chi_{\mathcal{V}}$ cannot express because they concern the carrier's
runtime representation rather than its field values.

\begin{corollary}[Uniqueness of the Ring-4 configuration]
\label{ch04:cor:unique}
The fiber of $\mathcal{V}$ over ring $4$ is a singleton:
$\mathcal{V} \cap (\{4\} \times S \times L \times D) =
\{(4, \text{\textsc{open}}, \text{\textsc{not licensed}},
\mathrm{false})\}$.
Consequently ``a valid Ring-4 boundary'' is a definite description: any
two valid Ring-4 carriers agree on every field, and the use-site
predicate needs no tolerance, range, or case analysis---it tests one
point.
\end{corollary}

The corollary is why the v0.1 discipline is cheap enough to run
everywhere. A predicate that tested membership in a fifteen-element set
with exceptions would invite caching and shortcuts; a predicate that
tests a single known point, field by field, costs a handful of
comparisons and admits no interesting optimization, so there is no
engineering pressure to skip it. Representational austerity purchases
procedural discipline.

\section{The recheck-not-trust invariant}\label{ch04:recheck}

\subsection{Statement}

\begin{invariant}[Recheck, do not trust]\label{ch04:inv:recheck}
Every public operation of the kernel derives the boundary it enforces from
its own inputs and rechecks that boundary at the point of use. No operation
trusts a boundary because some earlier operation checked it; no operation
consults a shared mutable object to learn a status; and no operation accepts
a recomputable self-hash as evidence that a boundary is still what it was at
construction.
\end{invariant}

The invariant has three clauses, and each is worth unpacking, because each
closes a distinct loophole through which construction-time validation alone
would leak.

\subsection{Derive and recheck at use}

Construction-time rejection (Section~\ref{ch04:consistency}) guarantees that
no invalid combination is ever a value. It does not guarantee that a value
checked yesterday is valid today, for two reasons. The first is mundane: the
host language provides no in-process memory protection. The carrier is a
frozen record, which resists \emph{accidental} mutation---ordinary attribute
assignment raises---but a determined caller in the same process can bypass
the freeze through the language's low-level attribute machinery. The kernel
does not pretend otherwise. The second reason is architectural: values flow.
A boundary object constructed in one component is embedded in an artifact,
passed across an interface, stored, retrieved, and eventually consulted by a
component that never saw its construction. If validity were checked only at
construction, the checking component and the relying component would be
different, and the reliance would be exactly the kind of unexamined trust
the kernel is built to avoid.

The v0.1 discipline is therefore that every use site re-runs the exact-field
predicate of Rule~\ref{ch04:rule:predicate}. A crosswalk validates its
boundary when constructed \emph{and} the registry validates the boundaries of
its records when loaded \emph{and} the registry re-validates them when a
record is resolved \emph{and} the record re-validates its own boundary inside
every integrity check. The pattern is deliberately redundant. Redundant
checking of cheap predicates is the kernel's substitute for the trusted
intermediary it refuses to have: instead of one blessed component whose
word is final, there are many unprivileged components each of which verifies
what it relies on. The cost is a few field comparisons per operation; the
benefit is that a corrupted boundary is caught at the first use after
corruption, with a typed error naming the artifact, rather than persisting
silently inside a structure that was ``already checked.''

\begin{remark}
The recheck discipline also changes what a successful attack would have to
look like. Against a check-once design, an attacker mutates a value after
its single validation and wins. Against a check-at-every-use design, the
attacker must either mutate the value into another \emph{valid}
configuration---but the only valid configurations reachable from Ring-4 open
change ring, status, or license fields that every v0.1 use site pins to
exact Ring-4 values, so the mutated object fails the very next predicate---
or subvert the predicate itself, which is a plain function over primitive
fields with no dispatch through the attacked value. The design cannot stop
an adversary who can rewrite the kernel's own code in memory; no in-process
scheme can. It arranges instead that anything short of that produces a loud,
typed failure. This is a containment posture, not a security claim, and the
Gate-4 receipts exercise only its fixture-level behavior.
\end{remark}

\subsection{No shared mutable trust}

The second clause forbids a pattern so convenient that it must be named to be
resisted: a module-level ``current status'' object that components consult,
and that privileged code updates when status changes. Such an object is
shared mutable trust. It has three defects. It is a single point of silent
failure: one mutation, however caused, changes the effective status of every
artifact that consults it. It inverts the evidence relation: artifacts no
longer carry their status; they borrow it from ambient state, so the same
artifact can be ``licensed'' at one moment and not at another with no change
to the artifact itself, and no record of why. And it is unauditable: a
consultation of ambient state leaves no trace in the artifact of what was
consulted or when.

The kernel's arrangement is the opposite on all three points. Status is a
field of the artifact, immutable, validated at construction, and rechecked at
use. Two artifacts never share a carrier. There is no ambient status to
consult and therefore no ambient status to corrupt. When Chapter 6's
registries serialize an artifact, the boundary goes with it, byte for byte,
into the digested evidence---so the status an artifact claimed is part of
what its digest commits to.

\subsection{No recomputable self-hash as trust}

The third clause addresses a subtler temptation, and it is the hinge to
Section~\ref{ch04:tamper}. Suppose the carrier computed a digest of its own
fields at construction and stored it alongside them, and suppose use sites
verified the stored digest instead of re-running the field predicate. This
looks like an integrity mechanism, and against \emph{accidental} corruption
it is one: a stray mutation of a status field would, with overwhelming
probability, disagree with the stored digest.

Against an adversary it is nothing at all. The digest is a public function
of the fields. Whoever can rewrite the fields can rerun the function and
rewrite the digest in the same act. A stored self-digest therefore cannot
distinguish ``unchanged since construction'' from ``changed by someone who
also updated the digest''---which is to say, it detects exactly the
adversaries who are not trying. The program has direct historical experience
of the resulting failure mode: self-hashed receipts whose hashes were
perfectly valid, because of course they were; the entity that fabricated the
receipt computed them.

The kernel accordingly assigns self-digests a strictly limited role. Where
they appear (Chapter 5's canonical encodings give every evidence record a
deterministic digest), they serve comparison and deduplication and the
detection of accidental divergence between independently held copies. They
are never consulted \emph{in place of} field-level revalidation, and they are
never, anywhere, treated as evidence of authorization. What separates a
digest that means something from a digest that does not is possession by a
party the mutator cannot rewrite---an external pin, in the vocabulary of
Chapter 6, where the registry loader receives its expected hashes from the
caller precisely so that the loaded bytes cannot vouch for themselves.

\section{Status change without mutation}\label{ch04:lifecycle}

A natural question about an immutable status carrier is how status ever
changes. An artifact that is \textsc{open} today may, one day, pass
review; the disclosure of a held result may one day be authorized. If
the carrier cannot be updated, where does the change live?

The kernel's answer is that it does not live in the kernel at all, and
this is a design position rather than a gap. The v0.1 kernel contains no
transition function on boundaries: no operation takes a boundary and
returns a successor boundary, no operation rewrites a boundary in place,
and no artifact's boundary field is assignable after construction. A
change of status, when the surrounding process one day produces one, is
an \emph{act}---a review, an admission decision, a disclosure
authorization---and acts are recorded the way the program records
everything else: as new artifacts. The reviewed successor of a Ring-4
candidate is a new artifact version carrying a new boundary, constructed
under Rules~\ref{ch04:rule:c1}--\ref{ch04:rule:c3} like any other, with
the admission act itself recorded as data in the ledgers of Chapter 6.
The predecessor persists, unchanged, with the status it always had.

Three properties fall out of this arrangement, each of which an in-place
transition model would forfeit.

\emph{History is retained by construction.} Because status change means
new-artifact-with-new-boundary, the trajectory of an artifact through
the admission process is a sequence of immutable records rather than the
final state of an overwritten field. An auditor reconstructing ``when
did this become reviewed, and what did it look like before?'' reads the
sequence; nothing must be recovered from logs of destroyed states. This
is the provenance instinct---keep the entity, the activity, and the
derivation distinct and durable \cite{prov}---applied to status.

\emph{No time-of-check races on status.} A mutable status field invites
the classic race: one component checks the boundary, another mutates it,
the first acts on stale information. With immutable carriers the checked
value cannot change between check and act; what can happen is only that
a \emph{newer artifact version} exists, which is a visible, queryable
fact rather than a silent invalidation of a value in hand.

\emph{No demotion laundering.} Transitions run both ways, and the
dangerous direction is not always upward. A program that can promote a
boundary in place can also quietly \emph{demote} one---or, more
insidiously, revert a demotion---leaving no trace that the elevated
status ever existed or was ever withdrawn. In the new-artifact model a
withdrawal is itself an act producing a record; the demoted
predecessor's history does not vanish because history is never stored
in a mutable cell.

\begin{remark}[The transitional states]\label{ch04:rem:transitional}
The admission-state vocabulary includes \textsc{in-review}, and
Proposition~\ref{ch04:prop:count} counts valid configurations containing
it at rings 1 through 3. In v0.1 those configurations are representable
but unconstructed: no kernel artifact carries them, because every v0.1
construction site pins the exact Ring-4 configuration. Their presence in
the vocabulary is forward provision, not dead code---a future admission
process needs a representable ``review in progress'' that is distinct
from both \textsc{open} and \textsc{accept}, and
Rule~\ref{ch04:rule:c3} already couples it correctly (no disclosure
during review). The kernel ships the vocabulary and the consistency
rules now, so that the future process inherits its state space from a
reviewed design instead of improvising one.
\end{remark}

\section{Tamper evidence versus authority}\label{ch04:tamper}

\subsection{Two questions}

The distinction this section develops is the load-bearing conceptual
distinction of the entire kernel, and the boundary object is where it lives.
It can be put as two questions that sound similar and share no evidence.

\begin{definition}[Tamper evidence]\label{ch04:def:tamper}
A mechanism provides \emph{tamper evidence} for an object if it answers the
question: \emph{has this object changed since a reference point?} The answer
is relative to a reference (a construction event, a pinned digest, a retained
byte string) and is available to anyone holding the reference.
\end{definition}

\begin{definition}[Authority]\label{ch04:def:authority}
A mechanism provides \emph{authority} for a status if it answers the
question: \emph{who conferred this status, and were they entitled to?} The
answer requires an authenticated link to a principal outside the object---a
reviewer, an institution, a legal act---and is unavailable to any mechanism
whose entire input is the object itself.
\end{definition}

The two questions are orthogonal in both directions. An object can be
perfectly tamper-evident and wholly unauthorized: the historical self-issued
receipts were bit-for-bit stable and hash-consistent from the moment of
fabrication, and fabricated. An object can be authorized and
tamper-corrupted: a genuinely reviewed record, later damaged in storage,
still \emph{was} authorized; what has been lost is the ability to
demonstrate which bits the authorization covered. Integrity mechanisms and
authorization mechanisms are different technologies with different trust
roots, and treating the first as the second is among the oldest failure
patterns in security engineering \cite{anderson}.

\subsection{What a self-digest answers}

Applied to the boundary carrier, the distinction yields a precise account of
what each mechanism in this chapter does and does not establish.

The consistency rules of Section~\ref{ch04:consistency} establish
\emph{coherence}: the recorded status is not self-contradictory. They answer
neither question. A coherent status can be both changed and unauthorized.

A digest of the carrier's canonical encoding, held by an external party or
pinned in an external record, establishes \emph{tamper evidence} relative to
that pin: if the object later hashes differently, it changed. It answers
Definition~\ref{ch04:def:tamper}'s question exactly and
Definition~\ref{ch04:def:authority}'s question not at all. Nothing in a hash
says who was entitled to create the hashed bytes. Content addressing is a
naming discipline, not an endorsement; the reproducible-builds literature is
built on precisely this separation---bit-identical artifacts establish that
a build was reproduced, while who ought to trust the build's inputs is a
separate, socially anchored question \cite{reprobuilds}.

A digest stored \emph{inside} the object itself establishes less: tamper
evidence against accidental mutation only, for the reasons of
Section~\ref{ch04:recheck}. The kernel's documentation is deliberately blunt
on this point, and the monograph repeats it as an invariant because
experience shows it is exactly the point that erodes.

\begin{invariant}[A self-digest is not authority]\label{ch04:inv:selfdigest}
A self-digest detects accidental or malicious mutation after construction
relative to an externally held reference; a self-held digest detects
accidental mutation only. In neither role is a digest a signature, a
reviewer identity, or an authority grant. No kernel operation treats hash
validity as authorization.
\end{invariant}

\subsection{What would answer the authority question}

It is worth stating positively what \emph{would} answer
Definition~\ref{ch04:def:authority}'s question, because the v0.1 kernel's
refusal to fake an answer is a design decision, not an oversight. An
authority mechanism requires a trust root: some principal whose statements
are accepted axiomatically, plus an authenticated channel from that principal
to the artifact---in practice a signature scheme with managed keys, an
enrollment process binding keys to reviewer identities, and revocation. It
further requires provenance structure: a record not merely that a status
holds but of the acts and agents that produced it, in the sense in which
provenance models distinguish an entity from the activity that generated it
and the agent that bears responsibility \cite{prov}. The v0.1 kernel has
none of this, on purpose. A half-built authority mechanism is worse than
none, because it manufactures artifacts that look like answers to the
authority question while answering only the tamper question. The kernel's
stance---developed as the constant-false authorization invariant in
Chapter 9---is that until a reviewed design supplies a real trust root,
the authority question has a standing answer: \emph{no}.

\begin{remark}
The two-question analysis also explains a choice that otherwise looks
strange: why the kernel bothers with elaborate integrity machinery at all if
integrity confers no authority. The answer is that the authority question is
answerable only \emph{about} specific bytes. A reviewer who one day signs a
claim record signs particular content; if the content is not fixed---if
parses can drift from bytes, encodings can vary by platform, statuses can
mutate in place---then even a genuine authorization act fails to bind
anything definite. Tamper evidence is not authority, but it is the substrate
on which any future authority must stand. Chapters 5 and 6 build that
substrate; this chapter's boundary object ensures that until authority
exists, the substrate cannot be mistaken for it.
\end{remark}

\section{How boundaries thread through the kernel}\label{ch04:threading}

The boundary object would be decoration if it were carried only by some
artifacts, or carried but not enforced. The v0.1 design threads it through
every public artifact type, with the same exact-field predicate at every
seam. This section walks the seams; the companion chapters treat each
artifact in depth.

\subsection{Semantic crosswalks}

A semantic crosswalk (Chapter 3) records a typed, non-identity relation
between references in different namespaces. Every crosswalk carries a
boundary, defaulted fresh to Ring-4 open, and the crosswalk constructor
rejects any boundary that is not exactly Ring-4 open. The consequence is
that no cross-namespace correspondence in v0.1 can so much as suggest a
reviewed status: even a correspondence the program's authors consider
obviously correct is constructible only as an open, unlicensed,
undisclosed research object. The historical Drive--Foundation conflation
could not have been laundered through a v0.1 crosswalk with an elevated
status, because no elevated status is constructible on that artifact type.

\subsection{Registry records and registries}

Chapter 6's claim and theorem registries enforce the boundary twice over, at
two different layers that must not be confused. The \emph{registry artifact}
---the loaded, validated, in-memory view---carries a kernel boundary, pinned
exactly to Ring-4 open like every other v0.1 artifact. The \emph{records
inside the ledger}, by contrast, carry their own status fields as data:
a claim record has a product-claim status that may spell \textsc{licensed},
an admission ring that may be 3, a review state that may be \textsc{reviewed}.
The loader validates these data fields against the record contracts---a
\textsc{licensed} claim row lacking reviewers, evidence packets, or a review
date is rejected as contractually malformed---but validation of the data is
not adoption of the data. The registry view of an all-\textsc{licensed}
ledger is still a Ring-4 open artifact. The kernel reads ledgers; it does
not become what it reads. This layering is the metadata-is-not-authority
invariant of Chapter 6, and it is the boundary object that makes the
layering representable: without a kernel-level boundary distinct from
record-level status fields, ``the registry says licensed'' and ``the kernel
asserts licensed'' would be one statement instead of two.

\subsection{Compatibility receipts}

Every compatibility execution under the deterministic-closure protocol
(Chapter 7) records the research boundary in its receipt, alongside the
compiler identities, worker identities, and result digests. The boundary is
recorded so that a receipt is self-describing about its claim status: a
consumer holding only the receipt can see---inside the digested content, not
in a detachable wrapper---that the execution it evidences was a Ring-4
research act. A receipt cannot be re-presented as evidence of licensed
behavior without altering its bytes, which the digest pin makes evident.
This is tamper evidence in service of honest labeling: the receipt does not
become more authoritative by carrying its boundary, but it becomes harder to
relabel.

\subsection{The runtime gate}

Chapter 9's decision model consumes boundaries at the last seam: an action
request carries typed intent, and the gate distinguishes nonclaim research
evidence from claim-bearing intent. There the boundary discipline reaches
its point. Structural validity of every artifact in the chain---references,
crosswalks, encodings, receipts, registry records---feeds the gate; and the
gate still answers \emph{no} to claim-bearing intent, constantly, because
the authority question has no affirmative answer in v0.1. The boundary
object is what lets the gate's refusal be principled rather than arbitrary:
the gate can point to the exact field, on the exact artifact, that records
the absence of authorization.

\subsection{Direction of dependence}

One structural property of the threading deserves note. The boundary type is
defined in the kernel's semantic layer and consumed by every other layer;
no layer above redefines or extends it (Rule~\ref{ch04:rule:exacttype}
forecloses extension by subclassing). Companion candidates---the PCBA board
model, the measurement stack, the topos layer---import the same carrier and
the same predicate. Status vocabulary is therefore global and singular: there
is exactly one way, program-wide, to say ``Ring-4 open,'' and it is checked
the same way everywhere. Vocabulary fragmentation---each subsystem with its
own status enum, related by convention---is how contradictory statuses
arise between subsystems without any single subsystem being internally
inconsistent.

\section{Worked misuse scenarios}\label{ch04:scenarios}

This section exercises the design against concrete misuse attempts. Each
scenario names the actor's move and the mechanism that blocks it. The
scenarios are fixture-shaped on purpose: each corresponds to a family of
tests in the bound Gate-4 suite, and the blocked-by clauses are statements
about tested constructor and predicate behavior, within the evidence bounds
of Chapter 12.

\begin{example}[Forged elevated status at construction]\label{ch04:ex:forge}
An agent constructs a boundary with ring 4 and product state
\textsc{licensed}, intending to attach it to a crosswalk so that downstream
tooling displays the crosswalk as licensed. \emph{Blocked at the first
step}: the carrier constructor raises on the R4/license family of
Table~\ref{ch04:tab:invalid}; no such value exists to attach. If the agent
instead constructs an internally consistent Ring-3 \textsc{accept}
\textsc{licensed} boundary, the value exists---and the crosswalk constructor
rejects it, because every v0.1 artifact constructor requires the exact
Ring-4 open boundary. The forgery fails at whichever seam it first touches.
\end{example}

\begin{example}[Impersonation by subclass]\label{ch04:ex:subclass}
An agent defines a subclass of the carrier that reports ring 4 and
\textsc{open} through overridden attribute access while storing elevated
values, or that overrides equality so that comparisons against the
distinguished Ring-4 boundary succeed. \emph{Blocked by}
Rule~\ref{ch04:rule:exacttype} at construction of the subclass instance
itself, and independently by Rule~\ref{ch04:rule:predicate} at every use
site, since the predicate's first conjunct tests exact runtime type and
never dispatches through the candidate value's own methods.
\end{example}

\begin{example}[Post-construction mutation]\label{ch04:ex:mutation}
An agent constructs a legitimate Ring-4 boundary, embeds it in an artifact,
and then uses low-level attribute machinery to overwrite the product state
in place, hoping the artifact's earlier validation carries forward.
\emph{Blocked by} Invariant~\ref{ch04:inv:recheck}: the next operation that
touches the artifact re-runs the exact-field predicate and raises a typed
boundary error. The mutation is not prevented---no in-process scheme
prevents it---but it cannot be \emph{used}: every path from the mutated
object to an effect passes through a recheck. If the agent also updates a
self-held digest to match, nothing improves for them:
Invariant~\ref{ch04:inv:selfdigest} says no kernel path consults the
self-digest in place of the field predicate.
\end{example}

\begin{example}[Ambient status injection]\label{ch04:ex:ambient}
An agent arranges for a shared ``process status'' object---a module global,
an environment flag, a configuration singleton---to indicate a licensed
context, expecting kernel operations to consult it. \emph{Blocked by}
absence of the seam: no kernel operation reads ambient status. The boundary
enforced by an operation is derived from the operation's own inputs
(Invariant~\ref{ch04:inv:recheck}, first clause). There is nothing to
inject into.
\end{example}

\begin{example}[Receipt laundering]\label{ch04:ex:laundering}
An agent presents a structurally valid artifact chain---green Gate-4 suite,
hash-consistent receipts, a coherent boundary---to a human or downstream
tool as evidence that the artifact is approved for product use, exploiting
the reader's inference from consistency to authorization. \emph{Blocked
within the kernel by} the constant-false authorization posture threaded
through Section~\ref{ch04:threading}: no kernel output states an
authorization, every receipt embeds its Ring-4 boundary in digested content,
and the runtime gate denies claim-bearing intent regardless of the chain's
structural quality. \emph{Not blocked outside the kernel}: a human reader
can still be talked into the inference. The kernel's contribution is that
every artifact the launderer can exhibit carries, in its own digested bytes,
the statement that it is open, unlicensed research---so the laundering
requires actively contradicting the artifact, not merely decorating it.
\end{example}

\begin{example}[Disclosure flip]\label{ch04:ex:disclosure}
An agent flips the disclosure flag on an otherwise-open boundary to make an
artifact appear publishable. \emph{Blocked by} Rule~\ref{ch04:rule:c3} at
construction (disclosure requires \textsc{accept}) and, for the mutation
variant, by the recheck discipline as in Example~\ref{ch04:ex:mutation}.
The subtle point is that the flag's semantics are conjunctive with status:
there is no representable state in which disclosure is authorized for an
artifact that has not passed admission, so ``leak first, regularize later''
has no target state to regularize into at Ring 4.
\end{example}

\begin{example}[Cross-artifact status borrowing]\label{ch04:ex:borrowing}
An agent pairs a Ring-4 artifact with a genuinely elevated artifact (say, a
Ring-3 record loaded from a ledger) and presents the pair so that the
elevated status appears to cover both. \emph{Blocked by} per-artifact
carriage: status is a field of each artifact, embedded in each artifact's
digested encoding; no kernel operation aggregates boundaries across
artifacts or computes a ``session status.'' The registry layering of
Section~\ref{ch04:threading} is the systematic form of this defense: even
the registry that \emph{contains} elevated records is itself Ring-4 open.
\end{example}

\begin{exercise}[Check-order absorption]\label{ch04:ex:absorption}
Table~\ref{ch04:tab:invalid} attributes each invalid combination to the
first rule that fires in the constructor's check order. Verify the
counts by re-deriving the table under the opposite order (license
precondition checked first, Ring-4 closure last), and identify the
combinations that migrate between families. Conclude that the
\emph{partition} of the 33 invalid combinations is an artifact of check
order while the invalid \emph{set} is not---and explain why fixtures
should therefore assert rejection, and may assert the error family only
where the check order is itself part of the documented contract.
\end{exercise}

\begin{checkclass}[Boundary fixture families]\label{ch04:check:fixtures}
The boundary behavior described in this chapter corresponds to the
following fixture families in the bound Gate-4 suite: (i) construction-
rejection fixtures covering representatives of each invalid family of
Table~\ref{ch04:tab:invalid}; (ii) exact-type fixtures rejecting
subclassed carriers and enumeration look-alikes (strings spelling
enumeration words); (iii) predicate fixtures exercising each conjunct of
the exact-field predicate independently; (iv) recheck fixtures mutating
a constructed carrier and asserting typed failure at next use; and (v)
seam fixtures asserting that each v0.1 artifact constructor (crosswalk,
registry record, registry) rejects every valid non-Ring-4 boundary.
Passing receipts for these families are historical software evidence
under the bounds of Chapter 12; they license no claim beyond tested
behavior on the fixtures.
\end{checkclass}

\begin{threat}[Residual: kernel-code subversion]\label{ch04:threat:residual}
An adversary who can modify the kernel's own code or the interpreter in
memory can disable the predicate, the constructors, and the gate together.
No mechanism in this chapter addresses that adversary, and no claim is made
otherwise. The design target is the historically observed failure class---
honest components and agents drifting into overclaim through convenient
representations---plus low-effort in-process tampering, not a hostile
runtime.
\end{threat}

\section{What the boundary object does not do}\label{ch04:limits}

The chapter closes by listing, without euphemism, the questions the boundary
object leaves open; they are drawn into the obligation set of Chapter 12.

First, it confers nothing. Fifteen valid configurations are representable;
in v0.1 exactly one is constructible on kernel artifacts, and none is
evidence of any status beyond the artifact's own assertion about itself.

Second, it authenticates nothing. There is no reviewer identity, no
signature, no trust root; the authority question of
Definition~\ref{ch04:def:authority} has no affirmative answer anywhere in
the kernel, and every mechanism that resembles an answer (digests, receipts,
coherent statuses) has been shown to answer a different question.

Third, it does not survive a hostile runtime
(Threat~\ref{ch04:threat:residual}), and its tamper evidence is only as good
as the external pins that hold the digests---pins whose discipline is the
subject of Chapters 5 and 6.

Fourth, its ring semantics are inherited, not defined: the kernel enforces
the consistency rules over the ring index but does not define what review at
each ring consists of. That definition lives in the Project Coherence
process documents, outside the kernel, and a future authorization design
(Chapter 9) must bind to it explicitly.

What the boundary object does do is narrow: it makes status explicit,
coherent, immutable-in-use, rechecked at every seam, and honest about the
difference between ``unchanged'' and ``authorized.'' The remaining chapters
build outward from that narrow guarantee: Chapter 5 fixes the bytes that
boundaries and everything else are committed to; Chapter 6 fixes the
ledgers that carry status vocabulary as data; and Chapter 9 shows the gate
that turns the whole discipline into a standing, principled \emph{no}.
